\documentclass[../main]{subfiles}
\begin{document}
\chapter{Motor asincr\'onico}
\begin{center}
\includegraphics[scale=0.4]{images/04/motor.jpg}
\end{center}
\info{Motor Trif\'asico}{Deslizamiento, principio de funcionamiento, arranque.}

\actividad{Trabajo pr\'actico{} de investigaci\'on:}
\begin{enumerate}
\item Explique qu\'e{} es el resbalamiento o deslizamiento en un motor asincr\'onico.
\item Investigar la secuencia paso a paso del campo magn\'etico{} giratorio en un motor trif\'asico.
\item Explicar la curva caracter\'istica{} de un motor trif\'asico{} asincr\'onico{} (curva par-velocidad).
\item Explicar el rendimiento de un motor y su influencia sobre el par.
\item Explicar c\'omo{} se calcula el par de un motor, sus unidades, y adem\'as{}:
\begin{itemize}
\item Par m\'aximo.
\item Par nominal.
\item Par de arranque.
\end{itemize}
\item Explicar la forma de conexi\'on{} de devanados en un motor trif\'asico{} y las formas de conexi\'on{} en la placa del motor.
\item Explicar los siguientes sistemas de arranque, realizar esquema de mando y potencia en caso de ser necesario.
\begin{itemize}
\item Arranque Directo.
\item Arranque estrella-tri\'angulo.
\item Arranque por devanados separados.
\item Arranque por variador de frecuencia.
\item Arranque suave.
\end{itemize}
\item Investigar informaci\'on{} de internet de al menos 3 motores trif\'asicos{}, obteniendo: resbalamiento, rendimiento, corrientes nominal, par m\'aximo, par nominal, par de arranque, corriente de arranque, y comparar en una tabla.
\item Investigaci\'on: se desea elevar una carga de $P=\qty{1000}{kg}$ a una velocidad lineal y constante de $v_0=\qty{10}{m/s}$. Se pueden usar mecanismos de reducci\'on{} y poleas. \textquestiondown{}Qu\'e{} motor deber\'ia{} comprar?
\item Si duplico el peso, \textquestiondown{}a qu\'e{} velocidad podr\'a{} ascender la carga?
\end{enumerate}
\end{document}